% Chapter 1: one-page quickstart.
\chapter{Quickstart}\label{ch:quickstart}
\index{quickstart}

\begin{tcolorbox}[enhanced,colback=white,colframe=black,boxrule=1.6pt,arc=0pt,
  fonttitle=\bfseries\Large,coltitle=white,colbacktitle=black,
  title={Wind one roll in seven steps}]
\begin{enumerate}[leftmargin=2em,itemsep=0.15em]
  \item \textbf{Power off. Load and thread the paper} past the measuring
        wheel the way you were shown at installation. Keep hands clear of
        the roll and the wheel.
  \item \textbf{Power on and wait one second.} The screen shows
        \shows{-= PAPER WINDER =-} and \shows{[1] 30yd [2] 90yd}.
        \led{1} and \led{3} are lit.
  \item \textbf{Push \btn{1} for 30 yards, or \btn{2} for 90 yards.}
        One push starts the wind. The motor runs and the screen shows
        \shows{WINDING...}, a percent bar and the length.
  \item \textbf{Let it finish.} About 2 yards before the end the motor
        stops, then the screen shows \shows{TRIMMING...} and the motor makes
        bumps that get shorter and shorter to finish the length (this can
        take up to about half a minute).
  \item \textbf{Three short beeps and \shows{DONE}} mean the roll is complete:
        at least the full length plus 1 inch went past the wheel.
        Cut the paper and remove the roll.
  \item \textbf{Long beeps and \shows{!! UNDER TARGET !!}} mean the roll is
        \emph{not} guaranteed. Push any button to silence the alarm, then
        see Section~\ref{sec:under-target}.
  \item \textbf{To stop at any time, push \btn{1} or \btn{2}.}
        In an emergency switch the machine's power off.
\end{enumerate}
\end{tcolorbox}

{\small\textbf{Other buttons.} \btn{3} pauses and resumes a wind. \btn{4} turns
the motor on and off by hand (``manual'') from the idle screen, or nudges
the paper while a wind is paused.

\textbf{One push = one action.} The panel buttons are latching switches: a
push leaves the button sitting in, the next push lets it pop out. The
machine counts \emph{every} push, in or out, as one click. Which way a
button is resting does not matter. After a wind it is normal for a button
to stay pushed in.

\textbf{You do not need the PC} to wind paper. The PC program is only for
calibration and settings (Chapter~\ref{ch:pc}).\par}
